\documentclass[10pt,a4paper]{article}
\usepackage[utf8]{inputenc}
\usepackage{amsmath}
\usepackage{amsfonts}
\usepackage{amssymb}
\usepackage{enumitem}
 \usepackage{amsthm}
 \usepackage{verbatim}
\usepackage{bbm}

\newtheorem{example}{Example}
\newtheorem{definition}{Definition}
\newtheorem{theorem}{Theorem}
\begin{document}



\section{Measure theory}
We recap some basics of measure and probabiity theory  following the structure of \cite{szepesvari2020bandit}, mathematical aid can be found in \cite{tao2011introduction} or watch \cite{brightsideofmaths}.
\begin{comment}

\begin{itemize}
	\item Borel / Lebeques $\sigma$- algebras
	\item Integration theory
	\item Fundamental Theorems (Carath\'eodory's extension theorem)
	\item Probability calculus (condtitionals )
	\item Random variables as components of random vectors.
	\item Transition probability kernel
\end{itemize}
 
\end{comment}

\subsection{Basic Defintions}
\begin{itemize}
\item \textbf{$\sigma$-algebra:} Given a set $X$, the $\sigma$-algebra $\mathcal F$ on $X$ is a collection of subsets of $X$, $2^X \subseteq \mathcal F$  such that 
	\begin{enumerate}
		\item $X\in \mathcal F$,
		\item It is closed under complementation:, $A \in \mathcal F \Rightarrow X \setminus A \in \mathcal F $,  
		\item It is closed under countable unions: $A_1,  A_2, A_3, \dots \in \mathcal F \Rightarrow A_1 \cup A_2, A_3,  \cup ...  \in \mathcal F $ , 
	\end{enumerate}
where $2^X$ is the power set of $X$. Apparently $\emptyset \in \mathcal{F}$ and  as a consequence also countable intersections belong to $\mathcal F$: $A_1,  A_2, A_3, \dots \in \mathcal F \Rightarrow A_1 \cap  A_2, \cap A_3, \cap...  \in \mathcal F$\footnote{For $A_1, A_2 \in \mathcal F \Rightarrow X \setminus A_1, X  \setminus A_2 \in \mathcal F   \Rightarrow (X \setminus A_1) \cup  (X \setminus A_2) \in \mathcal F \Rightarrow X \setminus ((X \setminus A_1) \cup  (X \setminus A_2)) = A_1 \cap A_2 \in \mathcal F$, where the last equality follows from of De Morgan's law.}.  Apparently $\{\emptyset, X\}$ is the smallest and $2^X$  the largest  $\sigma$-algebra of $\Omega$. 
\item The tuple $(X, \mathcal F)$ is called \textbf{measurable space}.
\item \textbf{Measure}. Let  $(X, \mathcal F)$  be a measurable space. A measure $\mu: \mathcal F \rightarrow [0, \infty]$ is a function satisfying
	\begin{enumerate}
		\item Non-negativity: $\mu(A) \geq 0 \quad \forall  A \in \mathcal F$
		\item Null empty set: $\mu(\emptyset) = 0 $
		\item Countable additivity ($\sigma$-additivity). If $\{A_i\}_{i=1}^\infty$ are \textit{disjoint} sets in $\mathcal F$, then 
		\begin{align} \nonumber
		\mu\!\left(\bigcup\limits_{i=1}^\infty A_i \right) = \sum_{i=1}^\infty \mu(A_i)
		\end{align}
	\end{enumerate}
A \textbf{probabilty measure} additionally satisfies
\begin{enumerate}[start=4]
	\item $\mu(X) = 1$
\end{enumerate}
\item The triple  $(X, \mathcal F, \mu)$ is called \textbf{measure space}.

\end{itemize}

\subsection{Relation between measurable spaces.}

\begin{itemize}
\item 
\textbf{Measureable function}. Let $(X, \mathcal F)$ and $(Y, \mathcal G)$ be two measureable spaces. A function  $f: X \rightarrow Y$ is called  measurable  if $   f^{-1}(B) \in \mathcal F, ~\forall B \in \mathcal G$. So a measureable function respect the $\sigma$-algebra structure (pull backs measurable sets in it's codomain to measurable sets in it's domain). One also says $\mathcal F-\mathcal G$ (or, $\mathcal F / \mathcal G$) measurable to stress the underlying $\sigma$-algebras.
\item 
\textbf{$\sigma$-subalgebra:} Let $(\Omega, \mathcal F)$ be a measurable space with $\sigma$-algebra $\mathcal F$ on $\Omega$. An other $\sigma$-algebra $\mathcal G$ on $\Omega$ is called $\sigma$-subalgebra of $\mathcal F$ if $\mathcal G \subseteq \mathcal F $. 

\item 
\textbf{Filtration.} Given a measurable space $(X, \mathcal F)$. A filtration is a family of $\sigma$-subalgebras $\left(\mathcal F_t\right)_{t\geq0}$ such that $\mathcal F_s\subseteq \mathcal F_t \subseteq \mathcal F$, for all $s \leq t$. 

In a \textit{stochastic process} this is used to interpret  $\mathcal F_t$ as the collection of events observable up to time $t$ and $\mathcal F_s\subseteq \mathcal F_t$ implies that information does not decrease over time (knowledge at time $s$ is also availble at time $t$). A stochastic process $\left(M_t\right)$ is called adapted to the filtration $\left(\mathcal F_t\right)$ if $M_t$ is measurable with respect to $\mathcal F_t$.

\item 
\textbf{Intersection of $\sigma$-algebras is a $\sigma$-algebra. }Let $\{ \mathcal F_i \}_{i \in I}$ be a family of  $\sigma$-algebras  of $X$ (countable or not).  Define the intersection as 
\begin{align*}	
	\bigcap\limits_{i \in I} \mathcal F_i = \left \{    A \subseteq X   |  A \in \mathcal F_i  ~\forall i \in I \right  \}
\end{align*}

\item
 \textbf{$\sigma$-algebra generated by set}. The \textit{smallest} $\sigma$-algebra generated by a set $\mathcal C \subseteq 2^X$ is the intersection of all $\sigma$-algebras containing $\mathcal C$. 
\begin{align*}
	\sigma(\mathcal C) = \bigcap \left \{ \mathcal F  ~| ~ \mathcal C \subseteq  \mathcal F,  \mathcal F \text{ is }       \sigma\text{-algebra}    \right\}
\end{align*} 
Note the $\mathcal C$ is a family of subsets of $X$. This is extremely useful for constructing $\sigma$-algebras for non-finite sets (and constructing measures). For example the power set of $\mathbb R$ is  a beast (still a $\sigma$-algebra) and it is much more convenient to work with the  $\sigma$-algebra generated by all open sets of $\mathbb R$, called the Borel $\sigma$-algebra $\mathcal B (\mathbb R)$.

\item 
\textbf{$\sigma$-algebra generated by a function. } Let $X$ be a set and $(Y, \mathcal G)$ a measureable space and a function $f: X \rightarrow Y$. Denote the set of preimages by $\mathcal C = \{f^{-1}(B): B \in \mathcal G \}$. The $\sigma$-algebra generated by $f$  on $X$ is  the $\sigma$-algebra generated by the set of preimages $\sigma(f) := \sigma(\mathcal C) $.  Apparently it is the smallest $\sigma$-algebra on $X$ that makes $f$ measurable.

\item 
\textbf{Product of $\sigma$-algebras} Let $(X_1, \mathcal F_1), ( X_2, \mathcal F_2)$ be measurable spaces, then the product sigma algebra is defined by
\begin{align*}
	\mathcal F_1 \otimes \mathcal F_2& := \sigma(\mathcal F_1 \times \mathcal F_2 ) \\
%	X_1\times X_2 &= \{(x_1, x_2)  ~|~ x_1 \in X_1, x_2 \in X_2 \},\\ 
	 \mathcal F_1 \times \mathcal F_2 &= \{A_1 \times A_2  ~|~ A_1 \in \mathcal F_1,  A_2 \in \mathcal F_2 \}, 
\end{align*}
That is the   $\sigma$-algebra generated by the set  $\mathcal F_1 \times \mathcal F_2$ on $X_1\times X_2 = \{(x_1, x_2)  ~|~ x_1 \in X_1, x_2 \in X_2 \}$\footnote{$\mathcal F_1 \times \mathcal F_2$ is not a $\sigma$ algebra. This can be seen from the counterexample $\mathcal F_1 =  \mathcal F_2 = 2^{\{1,2\}}$}. Thus the product of measurable spaces can be declared as $(X_1 \times X_2, \mathcal F_1 \otimes \mathcal F_2 )$ Because $(\mathcal F_1 \otimes \mathcal F_2)\otimes \mathcal F_3  = \mathcal F_1 \otimes \mathcal (F_2 \otimes \mathcal F_3)$ it is enough to write $\mathcal F_1 \otimes \mathcal F_2\otimes \mathcal F_3$. Furthermore, for Borel sets things are smooth, $\mathcal B(\mathbb R^{p+q}) = \mathcal B(\mathbb R^{p}) \otimes \mathcal B(\mathbb R^{q}) $.
	
\item 
\textbf{Pushforward measure.} Let $(X, \mathcal F, \mu)$ be a measure space,  $(Y, \mathcal G)$ a measurable space, and $f: X \rightarrow Y$ a measurable function. The pushforward measure of $\mu$ under $f$, denoted by $f_*\mu$ (or $\mu \circ f^{-1}$) is a measure on $(Y, \mathcal G)$ defined by $f_*\mu(B) = \mu f^{-1}(B), ~ \forall B \in \mathcal G$.\footnote{The pushforward measure is indeed a measure. Countable additivity  follows from $B_i \cap B_j  = \emptyset \Rightarrow f^{-1}(B_i) \cap f^{-1}(B_j) = \emptyset, ~B_i, B_j \in \mathcal G$, that is disjointness carries over to the preimages. To show this assume the contrary, $B_i \cap B_j  = \emptyset$ but  $f^{-1}(B_i) \cap f^{-1}(B_j) \neq \emptyset$. Thus there exists $\omega \in f^{-1}(B_i)$ but also $\omega \in f^{-1}(B_j)$, therefore   $f(\omega) \in B_i$ and $f(\omega) \in B_j$  and therefore $\omega \in B_i \cap B_j  \neq \emptyset$, which is a contradiction.}
\end{itemize}

\subsection{Probability theory}
\begin{itemize}
\item A measurable space is commonly written as $(\Omega, \mathcal F)$, where $\Omega$ is also called \textbf{sample space} (the set of all outcomes).   
\item 
A measure space $(\Omega, \mathcal F, p)$ with $p$ a probability measure is called  \textbf{probability space}. Probability measures are also called \textbf{probability distributions} 
For a probability measure the complement follows as\footnote{$\Omega = A \cup \Omega \setminus A \Rightarrow  p(\Omega) = \mu(A)  + p(\Omega \setminus A) \Rightarrow   p(\Omega \setminus A)  = 1- \mu(A)$ by property $p(\Omega) = 1$ of a probability measure.}  $p(\Omega \setminus A) = 1 - p(A)$. 

\item Given a probability space $(\Omega, \mathcal F, p)$,  a measurable  function $X$  is called 
\textbf{random variable} if $X: (\Omega, \mathcal F) \rightarrow (\mathbb R,  \mathcal B(\mathbb R))$, \textbf{random vector} if  $X: (\Omega, \mathcal F) \rightarrow (\mathbb R^k,  \mathcal B(\mathbb R^k))$ and \textbf{random element} $X: (\Omega, \mathcal F) \rightarrow (\mathcal X,  \mathcal G)$, where $(\mathcal X,  \mathcal G)$ is a measurable space.  So a random variable is a random element with the specification that the codomain is the  measurable space of the real line together with the Borel $\sigma$-algebra.  

Because X is measurable,  the pushforward measure $X_*p = p \circ X^{-1} := p_X$ can be canonically assigned to render the codomain a probability space in either case. The pushforward measure is also called \textbf{law }. 

Note, that capital letters are common in probability theory for functions, which encourages confusions. Furthermore, the canonical existence of a pushforward measure gives raise to further confusion ('probability distrubition of a random variable' instead probability distribution induced by a random variable). Furthermore, as all questions can be answered from $p_X$ alone this is often used as an excuse to not mention the underlying probability space   $(\Omega, \mathcal F, p)$ at all. 

Since random variables are functions, one can consider compositions, and algebraic operations on the real line of them. 
\end{itemize}


\begin{example}\textit{Coin toss.} \label{ex_cointoss}
Define the probability space $(\Omega, \mathcal F, p)$ by  
\begin{align*}
\Omega &= \{H, T\}, \\
 \mathcal F &=2^\Omega = \{ \{H\}, \{T\},  \{H, T\}, \emptyset \}, \\
 p(A) &=
 \begin{cases}
 	\frac{1}{2}, & A =  \{H\}, \{T\} \\
 	1, & A = \{H, T\}, \\
 	0, & A= \emptyset,
 \end{cases}
\end{align*}
where $A\in \mathcal F$. Define the random variable

\begin{align*}
	X : \Omega &\rightarrow \mathbb R   \\
	     \omega &\mapsto \begin{cases}
	     	1, & \omega = H \\
	     	0, & \omega = T 
	     \end{cases}
\end{align*}
We need to show that the function $X$ is a random variable, i.e., it is a measurable function to $(\mathbb R, \mathcal B(\mathbb R))$. Pick a  set $B \in \mathcal B(\mathbb R)$. Then $X^-1(B) = \emptyset $ if neither 1 or 2 are in $B$,    $\{H\}$ if $1 \in B$, $\{T\}$ if $0 \in B$, $\{H, T\}$ if both are in $B$. In either cases $X^-1(B) \in \mathcal F$. So indeed $X$ is a random variable.

For the pushword measure note that singletons are in $\mathcal  B$ and thus countable unions of it and thus (subsets of) the natural numbers. Thus $p_X(\{1\}) = p_X(\{2\}) =  1/2$. 
\end{example}


\begin{example}\textit{Rolling an even number with a die.} \label{ex_die}
Define the probability space $(\Omega, \mathcal F, p)$ by  
	\begin{align*}
		\Omega &= \{1,2,3,4,5,6\}, \\
		\mathcal F &= 2 ^\Omega \\
		p(A) &= \frac{|A|}{|\Omega|} = \frac{|A|}{6}, A \in \mathcal F
	\end{align*}
	Define the random variable
	\begin{align*}
		X : \Omega &\rightarrow \mathbb R   \\
		\omega &\mapsto \begin{cases}
			1, & \omega \text{ even} \\
			0, & \omega \text{ odd} 
		\end{cases}
	\end{align*}

By analogue arguments as in Example \ref{ex_cointoss}, $X$ is a measurable function and thus a random variable and the pushforward measure reads $p_X(\{1\}) = p_X(\{0\})$.
\end{example}
\noindent
In Example \ref{ex_die} we could also have used the $\sigma$-subalgebra $\mathcal F' = \{\emptyset, \Omega, \{1, 3, 5\}, \{2,4,6\}  \} $, with the same random variable (because $f^{-1}(1) = \{2,4,6 \}, f^{-1}(1) = \{1,3,5 \} $, both contained in $\mathcal F'$). 
\newline
\newline
\noindent
\textbf{Remark:}  random variable  (measurable functions) act on $X: \Omega \rightarrow \mathbb R$ while the probability distrubutions ((pushforward) measures) are defined on the corresponding $\sigma$-algebras ($\mathcal F, \mathcal B (\mathbb R) $) and thus on sets  
However, there is another notational perversion:   $p_X(1)$ instead of $p_X(\{1\})$ are common for measures. 
Note that the pushforward measures are the same in the above example, while the random processes are very different (coin toss vs die). This is quite powerful and motivates constructing general probability distributions on $\mathcal B (\mathbb R)$, i.e, the Binomial distribution in the above Example. 
Therefore in practice,  the probability space $(\Omega, \mathcal F, p)$ is hardly specified and one talks of 'a random variables which is Binomial distributed'. Apparently random variables don't make sense without specifying the underlying probability space. So, this is a rather cryptic wording for 1) there is a underlying probability space, 2) there is a random variable (measurable function) such that it's pushforward measure is the binomial distribution (on $\mathcal B(\mathbb R) $).


\begin{itemize}
	\item \textbf{Indicator function}. Given a set $\Omega$ and $A \subseteq \Omega$, the indicator function $\mathbbm{1}_A$ given by
	\begin{align*}
		\mathbbm{1}_A : \Omega &\rightarrow \{0,1\}   \\
		\omega &\mapsto \begin{cases}
			1, & \text{if } \omega \in A \\
			0, & \text{ otherwise} 
		\end{cases}
	\end{align*}
	The indicator function is a random variable on a measurable space $(\Omega, \mathcal F)$ $\Leftrightarrow$ $A \in \mathcal F$.
	\item For a \textbf{discrete} probability space $(\Omega, \mathcal F, P)$ the $\sigma$-algebra is usually the powerset $\mathcal F = 2^\Omega$. It is common to use capital $P$ for the probabiliy measure.    It is enough to specify $P$ on singletons ($\{\{ x\} | x \in \Omega \}$) as the rest follows from the definition of a probabiliy measure.  Brackets for singleton measures are often sloppily omitted, $P(x)$ instead of $P(\{x\})$. 
\end{itemize}



\section{Stochastic process}
		
	
	
	\begin{definition}[Stochastic Process]
		Let $(\Omega, \mathcal{F}, \mathbb{P})$ be a probability space and $(S, \mathcal{S})$ a measurable space.  
		A \emph{stochastic process}  $X=(X_t)_{t \in T} $ is collection of random variables (elements)
		\begin{align*}
		\{ X_t : \Omega \to S \}_{t \in T},
		\end{align*}
		indexed by a set $T$ (typically $\mathbb N_0$, or $\mathbb R _{+}$). $S$ is called state space.
	\end{definition}

\noindent	
Thus, a stochastic process is a map $X: \Omega, T \rightarrow S$, $(\omega, t) \mapsto X_t(\omega)$ such that $X_t(\omega)$ is a $\mathcal F$--$\mathcal S$  measurable function. For a fixed $\omega$ the map $X(\cdot, \omega): T \rightarrow S, t\mapsto X_t(w)$ is called \textbf{path} of a stochastic process.  A stochastic process is called  \textbf{stationary in the strict sense} if for any $t_1, \cdots, t_m \in T$ and translation 
	$ t_{1+s}, \cdots, t_{m+s}\in T$ the random variables   $X_1, \cdots, X_m $ and
$X_{1+s}, \cdots, X_{m+s}$	have the same probability distributions.
	
	
\begin{definition}[Probability Kernel \cite{szepesvari2020bandit}]\label{def_probability_kernel}
	Let $(\mathcal S_1, \bar{\mathcal S}_1)$ and $(\mathcal S_2, \bar{\mathcal S}_2)$ be measurable spaces. A probability kernel (or markov kernel) $k$ from $(\mathcal S_1, \bar{\mathcal S}_1)$ to $(\mathcal S_2, \bar{\mathcal S}_2)$ is a map $k:  \mathcal S_1 \times \bar{\mathcal S}_2 \rightarrow [0,1]$ such that
	\begin{enumerate}
		\item $k(x, \cdot )$ is a probability measure on $(\mathcal S_2, \bar{\mathcal S}_2)$  for every fixed  $x \in \mathcal S_1$
		\item $k(\cdot, A )$ is  $\bar{\mathcal S}_1$--$\mathcal B([0,1])$-measurable for every fixed $A \in \bar{\mathcal S}_2$
	\end{enumerate}
\end{definition}	

\noindent
Notationally, $k$ is also called $(\mathcal S_1, \bar{\mathcal S}_1) \rightarrow (\mathcal S_2, \bar{\mathcal S}_2)$ probability kernel (while operating only on $ \mathcal S_1 \times  \bar{\mathcal S}_2$ ).  First, clarify the notation:
\begin{align*}
	k(x, \cdot )&: \bar{\mathcal S}_2 \rightarrow [0,1],~ A \mapsto K(x, A)\\
	k(\cdot, A )&:  \mathcal S_1 \rightarrow [0,1], ~ x  \mapsto K(x, A). 
\end{align*}
The idea behind the first requirement is to describe stochastic transitions; having arrived at $x_1 \in \mathcal S_1$ the next state $x_2 \in \mathcal S_2$ is sampled from $\sim k(x_1, \cdot)$. The second requirement essentially assures integrability, which is desireable to 1) construct products or compostions of proabability kernels and 2) to extend the theory in case $(\mathcal S_1, \bar{\mathcal S}_1)$ is a probability space, i.e it is additionally equipped with probability measure (in other words: if we wish to make statements about transitions if the state $x_1$ is uncertain).  

If $(\mathcal S_1, \bar{\mathcal S}_1, p)$ is a measure (probability) space and   $k$ is a $(\mathcal S_1, \bar{\mathcal S}_1) \rightarrow (\mathcal S_2, \bar{\mathcal S}_2)$ probability kernel   then $p\otimes k$ is a (probability) measure on $(\mathcal S_1 \times \mathcal S_2,, \bar{\mathcal S}_1 \otimes  \bar{\mathcal S}_2)$ defined by 
\begin{align}
	p\otimes k(A) &= \int _{\mathcal S_1}\int_{\mathcal S_2} \mathbbm{1}_A((x_1,x_2)) k (x_1, dx_2) dp(x_1) ,  
\end{align}
where $A \in \bar{\mathcal S}_1 \otimes  \bar{\mathcal S}_2$. So this allows to construct a joint probabily measure from a probability kernel and an ('initial') probability measure $p$.


Turning to the \textbf{product kernel},  let  $k_1$ is a $(\mathcal S_1, \bar{\mathcal S}_1) \rightarrow (\mathcal S_2, \bar{\mathcal S}_2)$ probability kernel  and $k_2$ is a $(\mathcal S_2, \bar{\mathcal S}_2) \rightarrow (\mathcal S_3, \bar{\mathcal S}_3)$ probability kernel. Then  $k_1 \otimes k_2$ is a probability kernel $(\mathcal S_1, \bar{\mathcal S}_1) \rightarrow (\mathcal S_1 \times \mathcal S_2, \bar{\mathcal S}_2 \otimes \bar{\mathcal S}_2)$ defined by
\begin{align}
	k_1\otimes k_2(x,  A) = \int _{\mathcal S_2}\int_{\mathcal S_3} \mathbbm{1}_A((x_2,x_3)) k_2 (x_2, dx_3)  k_1 (x_1, dx_2) ,  
\end{align}
where $x \in \mathcal S_1 \times S_2 , A \in \bar{\mathcal S}_1 \otimes  \bar{\mathcal S}_2$.

\begin{theorem}[Ionescu-Tulcea]
	tbd
\end{theorem}

\begin{comment}
	
	\begin{definition}[Markov Chain]
		Let $(S, \mathcal{S})$ be a measurable space, $\nu$ a probability measure on $(S, \mathcal{S})$, and 
		\[
		P(\cdot \mid s): \mathcal{S} \to [0,1]
		\]
		a transition kernel.  
		A \emph{Markov chain} with initial distribution $\nu$ and transition kernel $P$ is a stochastic process $\{X_t\}_{t \in \mathbb{N}_0}$ on the path space $(S^{\mathbb{N}_0}, \mathcal{S}^{\otimes \mathbb{N}_0})$ such that for all $n \geq 0$ and $B_0, \dots, B_n \in \mathcal{S}$,
		\[
		\mathbb{P}(X_0 \in B_0, \dots, X_n \in B_n)
		= \int_{B_0} \nu(ds_0) \int_{B_1} P(ds_1 \mid s_0) \cdots \int_{B_n} P(ds_n \mid s_{n-1}).
		\]
		Equivalently, $\{X_t\}$ satisfies the \emph{Markov property}:
		\[
		\mathbb{P}(X_{t+1} \in B \mid X_0, \dots, X_t) 
		= \mathbb{P}(X_{t+1} \in B \mid X_t), 
		\quad \forall B \in \mathcal{S}.
		\]
	\end{definition}
	
\end{comment}





\section{Markov decision process}
\begin{comment}
\begin{itemize}
	\item  markov decision model 
	\item Markov decision process
	\item Graph theory
	\item Bellmans equations vs Hamilton Jacobi
\end{itemize}
\end{comment}
The following definitions are based on \cite{doringrf}
\begin{definition}[Markov decision model]
A Markov decision model is a tuple $(\mathcal S, \mathcal A, \mathcal R, p)$, with 
\begin{enumerate}
	\item 
		State space: $(\mathcal S , \bar{\mathcal S}) $ is a measurable space with $\sigma$-algebra $\bar{\mathcal S}$
	\item 
	Action space: $(\mathcal A , \bar{\mathcal A}) $ is a measurable space with $\sigma$-algebra $\bar{\mathcal A}$, which is constructed by requireing  for every state $s \in \mathcal S$ a measurable space  $(\mathcal A_s , \bar{\mathcal A_s}) $  called the action space of state $s$ and declare 
	\begin{align*}
	\mathcal A &= \bigcup_{s\in \mathcal S} \mathcal A_s \\ 
	\bar{\mathcal A} &= \sigma \left( \bigcup_{s\in \mathcal S} \bar{\mathcal A}_s \right)
	\end{align*}
	\item 
	Reward: $(\mathcal R , \bar{\mathcal R}) $ is a measurable space with $\sigma$-algebra  $\bar{\mathcal R} $, where $\mathcal R \subseteq  \mathbb R, 0 \in \mathcal R$ and it' restricted Borel-$\sigma$-algebra $\bar{\mathcal R}$.
	\item
	Reward function: $p$ is a  $(\mathcal S \times \mathcal A, \bar{\mathcal S} \otimes \bar{\mathcal A} ) \rightarrow (\mathcal S \times \mathcal R,  \bar{\mathcal S} \otimes \bar{\mathcal R})$ probability kernel
	\begin{align*}
		 p:  &\bar{\mathcal S} \otimes \bar{\mathcal R}   \times  (\mathcal S \times \mathcal A) \rightarrow [0,1], \\
		 &(B, (s, a) )   \mapsto p(B; s, a )
	\end{align*}
\textbf{	Todo: this definition of a probability kernel is just the way around as above (definition \ref{def_probability_kernel}) -- fix this!}
\end{enumerate}
\end{definition}

The key ingredient of the definition is the kernel $p$ with the following interpretation.
If the system is currently in state s and action a is played, then $p(\cdot ; s, a)$ is the joint
distribution of the next state $s'$ and the reward $r$ obtained. 
If all sets are discrete then $p(s', r; s, a)$ is the probability to obtain reward $r$ and go to state $s'$
if action a was taken in state s. According to the definition the reward can depend on the
current state/action pair and the next state, but in most examples the reward and the next state will be independent. 

A Markov decision model does not fully describe a
process that runs from state to state and returns rewards. Only the transitions $(s, a) \mapsto (s',  r)$ is described but not how the next action $a'$ is chosen. Therefore and additional ingredient needed, the policy that can be chosen by the actor

\begin{definition}[Policy]
	For a Markov decision model $(\mathcal S, \mathcal A, \mathcal R, p)$ a policy is 
	\begin{enumerate}
		\item An initial Markov kernel $\pi_0$ on $\bar{\mathcal A}  \times \mathcal S$
		\item A sequence of probability kernels $\pi = (\pi_t)_{t \in \mathbb N}$ on $\bar{\mathcal A}  \times ((\mathcal S \times \mathcal A)^{t-1} \times \mathcal S)$ such that 
		\begin{align*}
			\pi_t(\mathcal A_s; s_0, a_0, \cdots, s_{t-1}, a_{t-1}, s) = 1
		\end{align*}
		for all $(s_0, a_0, \cdots, s_{t-1}, a_{t-1}, s) \in (\mathcal S \times \mathcal A)^{t-1} \times \mathcal S$
		The set of all policies is denoted by $\Pi$.
	\end{enumerate}	
\end{definition}

\begin{definition}[Markov decision process]
\end{definition}

\bibliographystyle{plain} % You can choose a different style like 'unsrt', 'alpha', etc.
\bibliography{paper}

\end{document}